% OpenAI-inspired publication table template.
% This is an operational style template, not an official OpenAI asset.
%
% Required packages in the document preamble:
%   \usepackage{booktabs}
%   \usepackage{threeparttable}
%   \usepackage{siunitx}
%   \usepackage{tabularx}
%   \usepackage{array}
%
% Recommended global setup:
%   \sisetup{
%     detect-all,
%     table-number-alignment = center,
%     input-symbols = {\,\pm\text{--}},
%     group-separator = {,},
%     group-minimum-digits = 4,
%   }
%
% Do not use vertical rules. Do not use \resizebox as the first response to
% density. Shorten labels, remove false precision, transpose, or split first.

\begin{table}[t]
  \centering
  \begin{threeparttable}
    \caption{Claim-led or specific descriptive caption.}
    \label{tab:openai_style_example}
    \setlength{\tabcolsep}{5.0pt}
    \renewcommand{\arraystretch}{1.24}
    \begin{tabularx}{\linewidth}{
      @{}
      >{\raggedright\arraybackslash}X
      S[table-format=2.1]
      S[table-format=2.1]
      S[table-format=2.1]
      S[table-format=+2.1]
      @{}
    }
      \toprule
      \textbf{Benchmark} &
      {\textbf{Baseline}} &
      {\textbf{Prior}} &
      {\textbf{Focal}} &
      {\textbf{$\Delta$ vs. prior}} \\
      \midrule
      \multicolumn{5}{@{}l}{\textbf{Primary evaluations}} \\
      Task A & 71.2 & 74.6 & \bfseries 79.1 & +4.5 \\
      Task B & 63.0 & 68.4 & \bfseries 70.0 & +1.6 \\
      \addlinespace[3pt]
      \cmidrule(lr){1-5}
      \multicolumn{5}{@{}l}{\textbf{Robustness evaluations}} \\
      Task C & 54.8 & \bfseries 60.2 & 59.9 & -0.3 \\
      Task D & 82.0 & 83.1 & \bfseries 85.4 & +2.3 \\
      \bottomrule
    \end{tabularx}
    \begin{tablenotes}[flushleft]
      \footnotesize
      \item \textit{Note.} Values are actual scores in percent; higher is better.
      Bold marks the best value in each row. Deltas are subordinate and never
      replace the actual values. Use an em dash for unavailable values and
      define uncertainty, sample size, and evaluation conditions here.
    \end{tablenotes}
  \end{threeparttable}
\end{table}

% ---------------------------------------------------------------------------
% Compact uncertainty-aware alternative
% ---------------------------------------------------------------------------
%
% \begin{table}[t]
%   \centering
%   \begin{threeparttable}
%     \caption{Method performance with 95\% confidence intervals.}
%     \label{tab:interval_example}
%     \setlength{\tabcolsep}{5pt}
%     \renewcommand{\arraystretch}{1.24}
%     \begin{tabular}{@{}l S[table-format=2.1] @{\,\,$\pm$\,\,} S[table-format=1.1]@{}}
%       \toprule
%       \textbf{Method} & \multicolumn{2}{c}{\textbf{Accuracy (\%)}} \\
%       \midrule
%       Baseline & 72.1 & 1.4 \\
%       Prior    & 75.8 & 1.1 \\
%       Focal    & \bfseries 78.0 & 1.0 \\
%       \bottomrule
%     \end{tabular}
%     \begin{tablenotes}[flushleft]
%       \footnotesize
%       \item \textit{Note.} Mean $\pm$ half-width of the 95\% confidence interval
%       over five independent seeds. Higher is better.
%     \end{tablenotes}
%   \end{threeparttable}
% \end{table}
